% BOOK-FILE-ID: f-76d7340d229c
% 付録の最後：本書の解法の全体像

\section{本書の解法を見渡す}\label{節：本書の解法を見渡す}

本書では,\sectionref*{節：目指す形}で述べた「既知の形をつくる」という方針のもと,たくさんの漸化式を解いてきた.
最後に,それらの解法を一段高いところから眺めてみよう.
結論から述べれば,本書の解法のほとんどは,次の3段階でできている.
\begin{enumerate}
  \item うまい置き換えで,難しい操作を既知の操作に移す.
  \item 移した先で解く.
  \item 元に戻す.
\end{enumerate}
違いは「何を置き換えるか」だけであり,それは大きく3種類に分けられる.
1つの量を置き換えるか,添字ごとに置き換え方を変えるか,複数の量を組み替えるか,である.
順に見ていこう.

\subsection{1つの量を置き換える}\label{節：1つの量を置き換える}

\paragraphbf{共役という考え方}
$a_{n+1}=f(a_n)$の形の漸化式を考える.
ある置き換え$b_n=h(a_n)$によって,\,$b_n$が$b_{n+1}=g(b_n)$という既知の漸化式を満たしたとしよう.
一つの数列で確かめただけでは,ほかの値でも成り立つとは限らない.
そこで,考える範囲のすべての$x$で
\[
  h(f(x))=g(h(x))
\]
が成り立ち,さらに$h$の逆で元の値に戻せる場合を考える.
つまり,「$f$で進めてから$h$で置き換える」ことと「$h$で置き換えてから$g$で進める」ことが同じ結果になる.
このとき,\,$f$と$g$は\textbf{共役}であるという.
これは同じ操作の繰り返しを調べる力学系の言葉であり,複素数の共役とは別の意味である.
力学系では,置き換えとその逆が連続であることなども求める.
共役なら,\,$f$を何回繰り返しても,\,$h$で移して$g$を同じ回数繰り返し,\,$h$で戻せばよい.
$g$が既知の操作なら,一般項が求まるのだ.

\paragraphbf{本書の解法を並べる}
本編の解法を,それぞれの置き換えが定義され,元に戻せる範囲で並べてみよう.
特性方程式型では$p\ne1$,対数型では$p>0$かつ$p\ne1$とする.
一次分数型の$\alpha,\beta$は異なる2つの不動点,\,$\lambda$は差の比の公比である.
\begin{center}
  \renewcommand{\arraystretch}{1.6}
  \begin{tabular}{l|l|l}
    型&置き換え$h(x)$&移った先$g(y)$\\ \hline
    特性方程式型&$x-\alpha$&$py$\\
    逆数型&$\frac1x$&$\frac rpy+\frac qp$\\
    一次分数型&$\frac{x-\alpha}{x-\beta}$&$\lambda y$\\
    対数型&$\log_px$&$qy+1$\\
    ニュートン法&$\frac{x-\sqrt c}{x+\sqrt c}$&$y^2$
  \end{tabular}
\end{center}
三角関数型と基本対称式型では,向きが逆になる.
角度から値を作る$\varphi(\theta)=\cos\theta$や$\varphi(u)=u+\frac1u$を使うと,
\[
  f(\varphi(\theta))=\varphi(2\theta),\qquad f(\varphi(u))=\varphi(u^2)
\]
となる.\footnote{$\varphi$は2つの値を同じ値に移すので,厳密には半共役と呼ぶ.$\cos$の値から角度は1つに決まらないことを,本編でも注意した.}
いずれにせよ,難しい操作が,掛ける・2乗する・2倍するといった既知の操作に移っている.

\paragraphbf{置き換えは重ねられる}
表の右の列をよく見ると,まだ既知の形になりきっていないものがある.
逆数型と対数型の$g$は一次式なので,不動点があれば,もう一度その不動点を引いて等比型にできる.
不動点がなく係数が1の場合は,すでに等差型である.
ニュートン法の$y^2$も,正の範囲では対数を取ればただの2倍になる.
このように,表の多くの式は置き換えを重ねることで
\[
  y\mapsto y+d\quad(\text{等差型}),\qquad y\mapsto ry\quad(\text{等比型})
\]
のどちらかになる.
\sectionref*{節：各パターンの立ち位置}で見たように,基本パターンは式変形の行き先として働いている.

\paragraphbf{置き換えの代償}
ただし,置き換えにはそれぞれ使える条件があり,元に戻すための情報も要る.
\begin{itemize}
  \item 逆数型は$a_n\ne0$,対数型は$a_n>0$でなければ置き換えられない.
  \item 一次分数型の$\frac{x-\alpha}{x-\beta}$は,不動点が2つ異なるときにしか作れない.
  \item 三角関数型では,\,$\cos$の値から角度が1つに決まらないので,戻すときに選び方が要る.
  \item 差を取る置き換えでは定数の情報が消えるので,戻すときに初期値で補う.
\end{itemize}
置き換えて解けたら終わりではなく,戻れることまで確かめて初めて解けたことになるのである.

\paragraphbf{小話：どこまで置き換えられるか}
では,どんな$f$でも,うまい$h$を選べば等比型に移せるのだろうか.
答えは否である.例えば実数全体で考えると,\sectionref*{節：式の構造を使う}の$f(x)=x^2-1$は不動点を2つ持つ.
一方,\,$g(y)=ry$の不動点は$r\ne1$なら1つ,\,$r=1$ならすべての実数である.
可逆な置き換えは不動点を一対一に対応させるので,この2つは実数全体では共役にならない.
しかし,不動点の近くだけに限れば話は変わる.
大学以降の数学では,複素数の近傍で解析的な関数$f$について,不動点$\alpha$で$0<|f'(\alpha)|<1$なら,\,$f$は$\alpha$の近くで$y\mapsto f'(\alpha)y$という等比型と共役になることが知られている.\footnote{ケーニヒスの定理と呼ばれる.ここでは複素解析での形を述べた.}
\sectionref*{節：不動点付近の線形化}で,\,$f'(\alpha)$が不動点付近でずれが何倍になるかを表すと述べたことの,より強い形である.
さらに$f'(\alpha)=0$のときにも,条件を満たせば累乗に移せる場合がある.\footnote{複素解析でのベッチャーの定理では,例えば$f$が不動点$\alpha$の近くで解析的かつ非定数であり,\,$f(z)-\alpha$の最初の0でない項が$k$次($k\ge2$)なら,局所的に$y\mapsto y^k$と共役になる.}
ニュートン法の例では,実際に置き換えを計算することで$b_{n+1}=b_n^2$が得られた.

\subsection{添字ごとに置き換え方を変える}\label{節：添字ごとに置き換え方を変える}

$n$を含む漸化式では,置き換え方そのものを添字ごとに変える.
\begin{itemize}
  \item \textbf{階比型}：\,$h_{n+1}=f(n)h_n$となる$h_n$で割り,\,$b_n=\frac{a_n}{h_n}$を定数型に移す.
  \item \textbf{階差等比複合型}：不動関数$g(n)$を引いて,\,$b_n=a_n-g(n)$を等比型に移す.指数の付加項では,0でない$q^n$で割って特性方程式型や等差型に移す.
  \item \textbf{変数係数の一次型}：\,$h_{n+1}=p_nh_n$となる$h_n$で割り,\,$b_n=\frac{a_n}{h_n}$を階差型に移す.
\end{itemize}
\paragraphbf{正規化}\label{節：正規化}
変数係数の一次型$a_{n+1}=p_na_n+q_n$で,\,$h_{n+1}=p_nh_n$を満たす0でない$h_n$を用意し,両辺を$h_{n+1}$で割ると,
\[
  \frac{a_{n+1}}{h_{n+1}}=\frac{a_n}{h_n}+\frac{q_n}{h_{n+1}}
\]
となる.\,$b_n=\frac{a_n}{h_n}$と置けば,\,$b_n$の係数は1になり,階差型に移る.
このように,階比型の解で両辺を割り,\,$a_n$の係数を1にする操作を,本書では\textbf{\bookterm{normalize}{正規化}}と呼ぶ.
例えば$a_{n+1}=\frac{n+1}{n}a_n$なら,\,$h_n=n$で割って
\[
  \frac{a_{n+1}}{n+1}=\frac{a_n}{n}
\]
となり,\,$\frac{a_n}{n}$は一定になる.
割る数$h_n$の隣り合う比$\frac{h_{n+1}}{h_n}$が,元の係数$p_n$と一致することが手掛かりである.

割る場合は,掛けられる係数を$h_n$に肩代わりさせている.
引く場合は,足される付加項を$g(n)$に肩代わりさせている.
具体的には,\,$h_n$は付加項を除いた$h_{n+1}=p_nh_n$を満たす.
一方,\,$g(n)$は元の式全体$g(n+1)=p_ng(n)+q_n$を満たす.
そのため,\,$a_n-g(n)$では付加項が消えるのである.

総和型の$S_{n+1}-S_n=a_{n+1}$や,隣接2項の差を取る方法も同じ仲間である.
差を取るという置き換えで,和や定数を消している.
差を取ることと和を取ることが逆向きの操作であることは,\sectionref*{節：和と差から微積分へ}で見た.

\subsection{複数の量を組み替える}\label{節：複数の量を組み替える}

隣接3項型や連立型のように,複数の量が絡み合って進む場合は,それらを組み替えて,別々に等比型で進む量を作った.
その組み替え方は,\sectionref*{節：漸化式と行列}で見たように,行列の固有ベクトルを探すという1つの操作にまとまる.
本編の「特性方程式」が何を求めていたかで並べると,次のようになる.
\begin{itemize}
  \item \textbf{特性方程式型・一次分数型}：不動点の方程式.解は固有ベクトルの\textbf{成分比}であり,差の比の公比は固有値の比である.(\sectionref*{節：隣接2項間を行列で解く})
  \item \textbf{隣接3項型・隣接多項間型・連立型}：固有値の方程式.解は固有値そのもの,すなわち各向きが\textbf{何倍になるか}であり,本編で作った量や重みは左固有ベクトルである.(\sectionref*{節：隣接3項間を行列で解く},\,\sectionref*{節：連立漸化式を解く})
\end{itemize}
同じ「特性方程式」という名前でも,方向を求めるものと何倍になるかを求めるものがあったのだ.
1つの行列から連立型・隣接3項型・一次分数型を同時に読む例は,\sectionref*{節：1つの行列で見渡す}で確かめた.

固有ベクトルに分けると,行列を掛けることはそれぞれの方向への定数倍になる.
つまり,固有ベクトルが十分にあれば,複数の量を組み替えることも等比型に移すことになる.

ただし,いつも分けきれるわけではない.
隣接3項型の重解のように固有ベクトルが足りないときは,等比型に分けきれない部分が残る.
0でない重解$\alpha$の場合は,その累乗で割ってから差分を取ればよかった.\,$u_n=\frac{a_n}{\alpha^{n-1}}$と置くと$a_{n+1}-\alpha a_n=\alpha^n\Delta u_n$となり,残る部分は階差型の繰り返しになる.(\sectionref*{節：差分で線形漸化式を解く})
ここで扱った定数係数の線形漸化式には,\textbf{等比型に分ける方法}と\textbf{累乗で割って差分を取る方法}があり,特性方程式の解の重なり方が選ぶ手掛かりになる.
行き先は,前者なら等比型,後者なら階差型である.

こうして得られた各部分を足し合わせたものが一般項になる.
付加項のない線形の漸化式では,解を定数倍して足しても解になるからである.
このことは\sectionref*{節：答えの形を見通す}で,解く前に答えの形を見通す道具として使った.

\subsection{置き換えられないとき}\label{節：置き換えられないとき}

もちろん,すべての漸化式が既知の形に移せるわけではない.
そのときは,一般項以外のものを調べることになる.
\begin{itemize}
  \item 変わらない量や周期を探す.(\sectionref*{節：変わらない量を使う},\,\sectionref*{節：非線形の周期性})
  \item 不動点の近くでの振る舞いや,極限を調べる.(\sectionref*{節：不動点付近の線形化})
  \item 整数性や余りを調べる.(\sectionref*{節：ガウス記号を含む式})
  \item 数列全体を母関数にまとめ,別の方程式に移す.(\sectionref*{節：数列をまとめて見る})
\end{itemize}
最後の母関数は,数列全体をまとめて置き換える方法と見ることもできる.
総和を含む非線形の漸化式が階比型に帰着したように,置き換える対象を広げれば,扱える漸化式も広がるのだ.

\paragraphbf{まとめ}
本書の「既知の形をつくる」という方針は,共役や固有ベクトルへの分解という考え方とも結び付く.分解しきれないときには,差分を使える場合もある.
多くの解法では,その行き先として最初に学んだ基本パターンが現れる.
基本パターンは,新しい漸化式を解く際にも使うため,本書の最初に置いたのである.
